\section{Introduction}
\label{sec:introduction}

Agentic workflows allow complex tasks to be carried from a natural language prompt. This has
broad applicability to allow non-expert users to perform difficult work. To facilitate the
agents' ability to do things beyond text output, tools need to be provided in the form of an
API. We examine two ways to expose this interface: A model context protocol (MCP) server which
shows all available processes along with metadata, and a vector database populated with the
available processes and accessed using retrieval augmented generation (RAG). These two
approaches are compared within an agentic system to create a workflow for a geospatial analysis
task, and the input prompt is paraphrased in various ways to improve the robustness of the trial.

Paraphrases are difficult to formally define but can still be categorized. For example, a
sentence can swap words for synonyms or change the order of clauses. Defining different types
of paraphrases allows for the identification and creation of paraphrases via a systematic
method \citep{vila2014}. We use six such categories and their subtypes to programmatically
generate paraphrases. Paraphrases have been shown to have a noticeable effect on LLM
performance \citep{wahle2024}. By testing many paraphrases, we aim to avoid any influence
related to semantics on the outcome of the trial.
